\begin{proof}

\begin{enumerate}[label={$\langle1\rangle$\arabic*.},leftmargin=*]
\item We use precisely the relative Albanese construction of
\cite[Proposition~2, pp.~605--608]{Kawai1969}, and not its subsequent
corollary. Applied to bielliptic fibres, which have
$p_g=0$ and $h^0(\Omega^1)=1$, it gives a compact complex surface
$B$, a holomorphic map $h:B\to C$, and a meromorphic map
$\Phi:Y\dashrightarrow B$ over $C$. On the complement $C^\circ$
of finitely many values, $B^\circ\to C^\circ$ is a smooth family
of elliptic curves and
\[
 \Phi|_{F_c}:F_c\longrightarrow B_c
\]
is the Albanese morphism of the bielliptic surface, with connected
elliptic fibres. To specify the comparison being imported, the
construction forms the relative period-lattice quotient, glues its
local Albanese maps by fibre translations, and compactifies at each
missing value by a normalized finite base change and the projective
space associated with the direct images of a divisor which is the
inverse image of a local Albanese section. In making this last
step equivariant, use the sum of the finite orbit of that divisor;
a selected section need not itself be invariant under the covering
group. Resolve the normalized base change first if needed, so the
divisor defines a line bundle. The orbit sum restricts on every
general fibre to the pullback of a positive-degree divisor on its
Albanese elliptic curve; sufficiently large powers give the same
Albanese image. Its naturally linearized line bundle and direct
images carry the finite covering-group action, and their quotient
glues the construction to $B^\circ$. These are the maps asserted by
Proposition~2, with that equivariant choice made explicit. The
later assertion in Kawai's corollary about multiple fibres is not
used here.

\item Resolve $B$ over $C\setminus C^\circ$ and resolve the induced
meromorphic map by a sequence of blowups with smooth centres
supported over those same values. This produces a diagram
\[
 \widetilde Y\xrightarrow{\ g\ }\widetilde B,
 \qquad \pi:\widetilde Y\longrightarrow Y,
 \qquad (h\circ\rho)\circ g=f\circ\pi,
\]
where $\rho:\widetilde B\to B$ is a resolution and both spaces
are smooth and compact. The map $g$ has connected fibres:
its general fibres are the connected Albanese fibres just described,
and its Stein factorization has finite degree one over the normal
surface $\widetilde B$. Hence
$H^1(\widetilde B;\Q)\hookrightarrow H^1(\widetilde Y;\Q)$.
Blowing up smooth centres does not change first cohomology, so
$b_1(\widetilde B)=0$. The compact-surface K\"ahler criterion
\cite[Theorem~1.1]{Peters1990} therefore gives a K\"ahler form
$\omega_B$ on $\widetilde B$.

\item The blowup cohomology formula gives
\[
 H^2(\widetilde Y;\R)
 =\pi^*H^2(Y;\R)\oplus\bigoplus_j\R[E_j]
 =\bigoplus_j\R[E_j],
\]
where the $E_j$ are the exceptional prime-divisor classes (with
the corresponding total transforms for successive blowups).
Every $E_j$ is supported above one of the removed values in $C$.
Consequently its restriction to a general smooth fibre $F_c$,
identified with its unchanged inverse image in $\widetilde Y$,
is zero. In particular
\[
 [g^*\omega_B]|_{F_c}=0\quad\hbox{in }H^2(F_c;\R).
\]
Choose a K\"ahler form $\omega_F$ on the bielliptic surface $F_c$.
The Albanese morphism $g|_{F_c}$ is a submersion to the elliptic
curve $\widetilde B_c$. Thus its pullback of $\omega_B$ is a
nonzero semipositive $(1,1)$-form and
\[
 \int_{F_c}g^*\omega_B\wedge\omega_F>0.
\]
The vanishing cohomology class makes the same integral zero by
Stokes' theorem, a contradiction. All maps in this argument are
the displayed relative Albanese map and its modifications; no
comparison between the two elliptic fibrations of a bielliptic
surface is needed.

\end{enumerate}
\end{proof}
